\section{Evaluation Evidence}\label{chapter:eval}

\subsection{Questions and evidence classes}\label{sec:eval:sample}

We separate four questions. \emph{Legality} asks whether an illegal
institutional change is rejected and a legal change commits atomically.
\emph{Liveness} asks whether proposal, attack, adjudication, transition,
repair, and replay paths make progress. \emph{Detection performance} asks
how often known faults are detected and clean cases are falsely rejected.
\emph{Effectiveness} asks whether governed co-evolution improves research
quality, calibration, or cost. Tests principally address legality, and
fixed fault cases exercise particular liveness paths. Detection
performance and effectiveness require labelled cases, independent runs,
and controls.

\begin{table}[t]
  \centering
  \small
  \begin{tabularx}{\columnwidth}{@{}lXX@{}}
    \toprule
    \textbf{Evidence} & \textbf{Supports} & \textbf{Does not support} \\
    \midrule
    1,124 relevant tests & implemented rules on tested inputs &
      safety on untested inputs or research quality \\
    Five faults, one clean case & activation of five specified paths and
      passage of the clean case &
      false-positive or false-negative rates, unknown-fault coverage \\
    Exploratory traces & execution of some long paths &
      reproducibility, distribution-wide liveness, superiority \\
    Unrun controlled study & nothing &
      quality, calibration, cost, or component-level causal effects \\
    \bottomrule
  \end{tabularx}
  \caption{Claims supported and not supported by the current evidence.}
  \label{tab:eval-status}
\end{table}

\subsection{Lifecycle demonstrations}\label{sec:eval:construction}

The exact focused selector recorded in the artifact instructions collects
1,124 tests directly concerning RQGM and the final manuscript gate; all
passed in this working tree. The complete suite selector collected 4,805
tests: 4,787 passed, 16 were skipped, and two were expected failures.
The focused set partitions without overlap into claim gate 64;
kernel, transition, epoch, state, founding, and mode 285; adversarial and
governance 126; repair and clean room 66; utility evolution 85; evaluation
fixtures and metrics 157; prompt evolution 93; paper roles 94; and context,
budget, meta-evolution, proposals, and adapter 154. These groups sum to
1,124; the machine-readable artifact lists every test identifier.
\Cref{sec:appendix:specification} maps the mechanisms to the component
lifecycle table, authority matrix, and recovery procedure.

The previously cited full SHA does not resolve from the public repository
and does not identify this uncommitted revision. We therefore withdraw the
public-artifact claim. This working tree now contains a local,
content-addressed reproduction bundle with source hashes, dependency
inventory, exact selectors, every collected identifier, a fully expanded
machine-readable authority matrix, and focused/full logs. It identifies a
dirty tree by its aggregate content digest rather than
pretending the current Git SHA names it. It is not an immutable public tag,
DOI archive, container, or CI record, so third-party retrieval is not yet
supported. A release must publish and archive the bundle and prompt material
before the paper cites a public artifact identifier.

The two expected failures are unrelated architecture debts: one command-layer
module still imports the visualisation layer, and the evaluator still calls
its model-provider library directly rather than through the shared model
client. They do not bypass an RQGM kernel check, but they remain visible in
the full log and are not counted as passes.

Test count is evidence of implementation scale and regression protection,
not scientific performance. The present environment has no coverage
instrumentation, so line coverage, branch coverage, and mutation score are
unknown. Nor does the aggregate establish that every invariant has at
least one test. We therefore do not use total test count as a safety
proxy.

Founding counts depend on configuration. The base research configuration
registers 31 prompts, one utility policy, and 20 components. Enabling the
paper archive adds one prompt and component each for self-preference
prosecution, paper review, and paper writing, producing totals of 35
prompt-or-policy records and 23 components. The base component count is 20,
including the registered research generator; the earlier ``32 and 19''
figures used the pre-generator component count and have been corrected.

Earlier exploratory executions observed boundary transactions, utility
policy replacement, re-scoring from saved raw axes, and paper-role
successor adoption. This report does not bundle the complete checkpoints,
exact provider-side model versions, decoding settings, seeds, independent
run counts, task descriptions, or costs. We therefore treat those traces
as implementation-path observations, not reproducible experimental
results.

\subsection{Fault injection and construction guarantees}\label{sec:eval:repro}

Fixed labelled fixtures cover five application fault paths: metric manipulation that
cannot be recomputed from run values, an unsupported overclaim, a prompt
candidate containing sealed material, forbidden clean-room input, and a
record dependent on a retired prompt leaking into the candidate set.
Each was rejected through its expected numeric-evidence, candidate,
kernel, or impact-exclusion path. One clean numeric claim was not falsely
rejected.

Additional integrity negatives mutate the event type, identifier,
transaction identifier, payload, and predecessor digest; reorder committed
events; and mismatch transaction membership. Each is rejected. A T16 test
also verifies a forced epoch close/open with a fresh fingerprint, and the
machine-readable authority export is checked against every sparse
allow-list entry. These are deterministic regression cases, not event
fuzzing or model checking.

This establishes only 5/5 detections and 1/1 clean passage for those fixed
fixtures. The sample is too small and contains only one difficulty level
per fault, so we do not report false-positive or false-negative estimates.
Fabricated prior art, denial of the adjudication path, incorrect
responsibility binding, and semantically preserved clean-room
contamination do not yet have completed labelled detection experiments.

Construction guarantees also have boundaries. Fixed code rejects
registry mutation outside the transition engine, unlisted transitions,
candidate authority expansion, penalties from unadjudicated attacks, and
replay of uncommitted transactions. This is kernel determinism over saved
inputs, not regeneration of model output. The schema-v2 audit chain binds
every replay-relevant envelope field and detects the tested interior
modifications, but without an external anchor it cannot detect rollback of
the whole checkpoint to an old valid prefix. T16 now force-closes the epoch
and opens a freshly fingerprinted epoch in the same transaction.

\subsection{Limitations}\label{sec:eval:limits}

An acceptance-grade study must preserve a preregistered task panel, exact
model and provider version, temperature, sampling method, maximum output,
tool versions, seeds, independent run count, epoch and node counts,
candidate counts, attack/motion/judgment/retirement counts, and token,
time, and monetary costs. Each rate needs a denominator, run-to-run
variance, and confidence interval.

We implemented five paper-comparison conditions aligned with the original
RQGM study. P0 evolves the writer against a fixed critic; P1 adds critic
replacement only; P2 co-evolves both roles but retains utility rows scored
by displaced critics; P3 makes those stale rows ineligible, realizing full
RQGM; and P4 adds blocking Constitutional-kernel enforcement to P3. All
five are evaluation presets over the same paper archive, budget, seeds,
and model settings, not new production modes. The headline comparison is
P0/P3/P4, with P1/P2 as mechanism ablations. Paper writing and fixed
post-hoc rubric review can use separately pinned models; the intended
study assigns Claude Code to writing and Codex to a panel disjoint from
the co-evolving reviewer. The paper does not include completed multi-seed
real-model runs, so this is an implemented protocol rather than a result.
The mechanism arms follow the original study, but the shipped budget of
eight epochs and at most 12 expansions per epoch does not reproduce its
12,288-evaluation compute scale \cite{iacob2026red}. A publication run
must preregister a common evaluator-call budget and report realized calls
per arm.

Further component ablations should remove replay, responsibility binding,
clean-room replacement, mutable utility policy, anchors, and the fixed
manuscript gate one at a time. Responsibility attribution and impact
repair additionally need labelled synthetic causal graphs to measure
target accuracy, reach precision and recall, and over-removal.

The warranted conclusion is consequently limited. The expected checks
fire on five fixed application fault cases, the event-integrity negatives
are rejected, and all 1,124 relevant tests pass.
Improved research quality, general detection performance, long-run
liveness, and cost effectiveness remain unestablished.
